\section{Categorical models and exact tensor reconstruction}\label{sec:positive}
Positive records are established by tensor matches to constructions,
not by lack of an obstruction. We use Deligne products of the standard
rank-two and rank-three unitary categories, the $TT_3$ category of
\cite[Definition~2.5]{EM}, and the positive quadratic constructions
of \cite[Proposition~5.1, Corollary~5.8]{MY}. We use the latter results
only for existence, not as a negative classification oracle.
Every displayed tensor is reconstructed from its defining rules and
compared under unit-preserving permutations with the census.

\subsection{A central fermion and a cyclic-four quadratic ring}
Let $\mathcal D$ be the unitary near-group category
$\operatorname{Ad}(E_6)$ with basis $1,h,X$ and rules
\begin{equation}\label{eq:E6}
 h^2=1,\qquad hX=Xh=X,\qquad X^2=1+h+2X.
\end{equation}
Its noninvertible dimension is $d=1+\sqrt3$, and
$D=\dim\mathcal D=4+2d$.

\begin{lemma}[Central fermion]\label{lem:fermion}
There is an invertible $z\in Z(\mathcal D)$ with
$F(z)=h$, $z^2=1$, and $\theta_z=-1$. The full subcategory generated
by $z$ is the unitary symmetric category $\mathrm{sVec}$.
\end{lemma}
\begin{proof}
The formal codegrees are $D$, $D'=6-2\sqrt3$, and $2$.
Thus $I(1)=1\oplus A\oplus B$, with dimensions
$1,1+d,2+d$. Since $FI(1)=3\cdot1+h+2X$, comparison of the unit
multiplicities and of the coefficients of $1,d$ gives
\[
 F(A)=1+X,\qquad F(B)=1+h+X.
\]
All three low twists are one.

The induction Gram matrix in the basis $(1,h,X)$ is
\[
 \begin{pmatrix}3&1&2\\1&3&2\\2&2&8\end{pmatrix}.
\]
Only $B$ among the low objects occurs in $I(h)$, once. The remaining
endomorphism dimension is $3-1=2$, so the remainder consists of two
distinct simples, both with multiplicity one. Its restriction is
$FI(h)-F(B)=2h+X$. Consequently these two objects are $z,W$ with
$F(z)=h$ and $F(W)=h+X$. The object $z$ has dimension one and is
invertible. Its square restricts to the unit; the only central simple
restricting to the unit is the unit itself, as the decomposition of
$I(1)$ shows. Therefore $z^2=1$.

The pointed associator class on $\langle h\rangle$ is trivial: the
rank-one module category on $X$ supplies a two-cochain whose coboundary
is the restricted associator. In the positive spherical structure,
$\nu_2(h)=1$. Equation~\eqref{eq:indicator} at $k=2$ implies that every
twist in $I(h)$ squares to one, because all dimensions are positive and
their total is $D$. At $k=1$ it gives
\[
 0=(2+d)+\theta_z+(1+d)\theta_W.
\]
The known positive summand has exactly the total dimension of the other
two. Equality in the triangle inequality forces
$\theta_z=\theta_W=-1$. Finally an order-two invertible in a unitary
ribbon category with twist $-1$ has self-braiding $-1$ in its canonical
positive structure. Its generated pointed braided category is
$\mathrm{sVec}$.
\end{proof}

\begin{proposition}[Fermionic zesting]\label{prop:zest}
The ring $\RR{2}{2}$ has a unitary categorification. More generally, if
a unitary fusion category $\mathcal D$ has a central fermion $z$ of
order two, the rules
\begin{equation}\label{eq:zest}
 (U,a)\star(V,b)=(U\otimes V\otimes F(z)^{ab},a+b),
 \qquad a,b\in\ZZ/2,
\end{equation}
on $\mathcal D\boxtimes\Vec_{\ZZ/2}$ admit a unitary associator.
\end{proposition}
\begin{proof}
Use the central two-cocycle $\lambda(a,b)=z^{ab}$ and normalized
associativity scalar
$\tau(a,b,c)=i^{abc}$, where the exponents use representatives
$0,1$. The only central interchange obstruction is
$c_{\lambda(a,b),\lambda(c,d)}=(-1)^{abcd}$. A direct calculation gives
\begin{align*}
 &bcd+a(b\mathbin\oplus c)d+abc
 -(a\mathbin\oplus b)cd-ab(c\mathbin\oplus d)\\
 &\hspace{35mm}\equiv 2abcd\pmod4,
\end{align*}
where $\oplus$ is addition modulo two. Thus the coboundary of $\tau$
is precisely that interchange obstruction. The associative-zesting
construction of \cite[Definition~3.2, Proposition~3.4]{DGP}, in the
relative-center formulation also recalled in \cite[Definition~2.1]{AGM},
therefore supplies the pentagon and rigidity for \eqref{eq:zest}.
All structural maps used are unitary, including the scalar $\tau$, so
the resulting fusion category is unitary. This use of a known
construction does not claim to have independently solved its individual
$F$-symbol equations.

Apply Lemma~\ref{lem:fermion}. Let $g=(1,1)$, identify
$g^2=(h,0)$, and write $Y=(X,0)$ and $X_1=(X,1)$.
Then the invertibles form $C_4$,
\[
 gY=X_1,\quad gX_1=Y,\quad
 Y^2=X_1^2=1+g^2+2Y,\quad
 X_1Y=g+g^3+2X_1.
\]
These rules reconstruct the full tensor that is based-isomorphic to
$\RR{2}{2}$. The supplement checks that equality and all sixteen scalar
coboundary identities exactly.
\end{proof}

The proposition is a positive construction, not an assertion that this
fusion category or its ring has not appeared in another presentation.
It is particularly useful here because a ring that survived the first
obstruction pass is thereby removed from the residual list for the
correct reason: existence rather than an attempted exclusion.

\subsection{Finite-group and group-theoretical models}
For $G=(C_3\times C_3)\rtimes C_4$, with the order-four rotation
$(x,y)\mapsto(-y,x)$, the character ring is reconstructed from four
linear characters and two induced characters. Exact character inner
products give all fusion coefficients and identify $\RR{2}{11}$.

We also compute the untwisted group-theoretical categories
$\mathcal C(G,1,H,1)$ for
\[
 (G,H)=(A_4,C_2),\quad(A_5,D_{10}),\quad(A_5,S_3).
\]
Here $D_{10}$ has order ten. The finite-group construction is standard;
we recall its explicit tensor calculation to distinguish it from a
mere comparison of dimension types. Realize the category as
$H$-bimodules in $\Vec_G$. Its simples are supported on double cosets
$HgH$ and indexed by irreducible representations of
$H\cap gHg^{-1}$. The tensor product is relative tensor product over
the group algebra of $H$. Equivalently it is convolution of the
corresponding equivariant vector bundles. The character of a convolution
fiber is computed by summing over the middle $H$-orbits and taking
invariants under the stabilizer; averaging the finite stabilizer action
gives its multiplicities. This is a finite exact character-inner-product
calculation. The code lists the groups as permutations, computes the
double cosets and stabilizers, and uses their exact character values.
It then checks the resulting complete fusion tensors.

The three outputs are respectively $\RR{2}{36}$,
$\RR{2}{105}$, and $\RR{3}{141}$. The last model is also discussed in
\cite[Section~6.5]{Integral}. Finite groups, ordinary unitary
representations, and relative tensor product over a finite-dimensional
$C^*$-algebra give unitary fusion categories in these untwisted cases.

\subsection{ADE modules, cyclic Izumi models, and a finite group}\label{sec:additionalmodels}
The following constructions add five realizations at multiplicity two
and one at multiplicity three. Each constructor builds its fusion tensor
without reading the census, then searches for a full based-ring match.

\begin{proposition}[ADE modules and their zesting]\label{prop:ADEmodels}
The rings $\RR{2}{75}$, $\RR{2}{14}$, $\RR{2}{128}$ and $\RR{2}{140}$
admit unitary categorifications. They are, respectively, the $E_6$ module
category over $SU(2)_{10}$, its fermionic grading zest, the $D_6$ module
category over $SU(2)_8$, and the local $D_{10}$ module category over
$SU(2)_{16}$.
\end{proposition}
\begin{proof}
Use the unitary $SU(2)_k$ category with simples $V_0,\ldots,V_k$ and
fusion coefficients
\[
 N_{ij}^l=1\quad\Longleftrightarrow\quad
 |i-j|\le l\le\min(i+j,2k-i-j),\quad i+j+l\equiv0\pmod2.
\]
The $E_6$ algebra is $V_0\oplus V_6$ at level ten. The $D$-type algebra
is $V_0\oplus V_k$ at levels divisible by four. Their existence,
separability, module categories and the ADE module graphs are the
Kirillov--Ostrik construction \cite[Theorem~6.1 and Section~7]{KO}.
These algebras and their module tensor products have the standard
unitary structures.

For $E_6$, label the graph vertices $0,\ldots,5$ and its edges
\[
 01,\quad12,\quad23,\quad34,\quad25,
\]
with unit vertex zero. Multiplication by the free image of $V_1$ is the
adjacency matrix $B$. The rows $(B^j)_{0,*}$, $0\le j\le5$, form an
invertible matrix $V$. Centrality of the free image of $V_1$ and the
cyclicity of the unit vector imply that all multiplication matrices are
uniquely determined by
$N_i=\sum_j(V^{-1})_{ij}B^j$. The exact constructor evaluates these
matrices, checks nonnegative integral entries and all tensor axioms,
and matches $\RR{2}{75}$.

The parity of the original $SU(2)$ labels descends to the module category;
in these coordinates it is $(0,1,0,1,0,1)$. The free image of $V_{10}$
is the degree-zero invertible at vertex four. The canonical central
free functor sends the order-two fermion $V_{10}$ to a central fermion
$g$ in the module category. Hence associative zesting may multiply the
product of two odd objects by $g$. Its obstruction is the fermionic
sign $(-1)^{abcd}$, trivialized by the scalar cochain $i^{abc}$, exactly
as in Proposition~\ref{prop:zest}. All sixteen cochain identities are
checked. This constructs a unitary fusion category and its tensor is
$\RR{2}{14}$. The construction uses the actual grading and central
fermion, not just a ring automorphism.

For the $D$ cases put $m=k/2$. The orbit modules are $O_j$,
$0\le j<m$, and the fixed object $V_m$ splits as $U_+\oplus U_-$.
The free image of $V_j$ is $O_{\min(j,k-j)}$ unless $j=m$.
Products of two orbit modules are obtained by applying this induction
rule to the displayed $SU(2)$ coefficients. For an orbit module times
$U_+\oplus U_-$ use the same rule; the difference $U_+-U_-$ has
eigenvalue zero for odd $j$ and $(-1)^{j/2}$ for even $j$. At the two
levels used here, $k\equiv0\pmod8$, the resolved objects are self-dual.
The fixed-point resolution formulas of \cite[Section~7, $D_{2m}$ fusion
table]{KO} give the remaining products. Equivalently, in the local
sector of this congruence class,
\[
 U_\pm^2=U_\pm+O_0+O_4+\cdots+O_{m-4},\qquad
 U_+U_-=O_2+O_6+\cdots+O_{m-2}.
\]
Together with Frobenius reciprocity these formulas are the constructor's
complete specification. At $k=8$ use all six modules; at $k=16$ use the
local modules $O_0,O_2,O_4,O_6,U_+,U_-$. Their tensors match
$\RR{2}{128}$ and $\RR{2}{140}$ respectively.
\end{proof}

\begin{proposition}[The Izumi cyclic-five model]\label{prop:Izumimodel}
The ring $\RR{2}{143}$ is the Grothendieck ring of the commutative even
part of the unitary Izumi $C_5$ subfactor.
\end{proposition}
\begin{proof}
The existence of the cyclic-five Izumi subfactor and the relevant even
part are established in the literature summarized in
\cite[Section~6.6]{GS}. Its commutative fusion rules are those in
\cite[Theorem~7]{EG}. To specify the full tensor used for the match,
write $1,a,b_1,b_2,c_1,c_2$, identify indices $i$ and $-i$ modulo five,
and put $b_0=c_0=0$, $\Sigma=a+b_1+b_2+c_1+c_2$. The products are
\begin{gather*}
 a^2=1+\Sigma,\qquad ab_i=\Sigma-b_i,\qquad ac_i=\Sigma+c_i,\qquad b_ic_j=\Sigma,\\
 b_ib_j=\delta_{ij}1+(1-\delta_{ij})a+\Sigma-a-b_{i+j}-b_{i-j},\\
 c_ic_j=\delta_{ij}1+(1+\delta_{ij})a+\Sigma-a+c_{i+j}+c_{i-j}.
\end{gather*}
All products commute and all simples are self-dual. The independent
constructor specializes these formulas at five and verifies all $216$
coefficients of the matched tensor. Existence is the cited subfactor
construction, not a consequence of the based-ring axioms alone.
\end{proof}

\begin{proposition}[An exact finite-group model]\label{prop:PSLmodel}
The ring $\RR{3}{47}$ is categorified by
$\Rep(GL(3,\FF_2))\simeq\Rep(PSL(2,7))$.
\end{proposition}
\begin{proof}
Generate the group as all invertible $3\times3$ binary matrices, with
actual matrix multiplication. It has $168$ elements and six conjugacy
classes of sizes $1,21,56,42,24,24$, ordered by element orders
$1,2,3,4,7,7$. Put $a=(-1+\sqrt{-7})/2$ and $b=\overline a$.
The candidate character table is
\[
\begin{array}{c|rrrrrr}
 &1&2&3&4&7&7\\\hline
1&1&1&1&1&1&1\\
3&3&-1&0&1&a&b\\
\overline3&3&-1&0&1&b&a\\
6&6&2&0&0&-1&-1\\
7&7&-1&1&-1&0&0\\
8&8&0&-1&0&1&1
\end{array}.
\]
We verify its representation-theoretic meaning, not just orthogonality.
The program computes the multiplication constants of the six class sums
in $\CC[GL(3,\FF_2)]$. For each row $\chi$ it checks that
$|C|\chi(C)/\chi(1)$ is a unital character of this class algebra.
The six characters are distinct and therefore exhaust its primitive
central characters. By semisimplicity they are the normalized central
characters of the actual irreducible representations. Row orthogonality
and positivity of the displayed degrees fix the normalization, so the
rows are the actual irreducible characters. Taking their product inner
products gives nonnegative integer structure constants, evaluated
exactly in $\QQ(\sqrt{-7})$. The resulting full tensor is based-isomorphic
to $\RR{3}{47}$. A finite-group representation category has its usual
unitary realization.
\end{proof}

\subsection{Categorical structures, extensions, and central quotients}\label{sec:literature-models}
The following determinations use categorical constructions, followed by
exact identification of their full fusion tensors. A based-ring
isomorphism transports both existence and a proved obstruction to a
specified categorical structure. It does not turn spherical existence
into pseudo-unitary existence. All identifications in this section are
recomputed by \filename{software/literature_models.py}; no census label
is an input to the construction.

\subsubsection{Schopieray's spherical model}
Schopieray constructs the rank-six modular category
$\mathcal C(\mathfrak{so}_5,3/2)_{\mathrm{ad}}$ in
\cite[Example~1, equations~(1)--(2)]{SchNonPU}. Its displayed fusion
matrices, in the order $(X_0,\ldots,X_5)$, give our tensor $\RR{2}{124}$
after the permutation
\[
 (0,1,5,4,2,3).
\]
Here, and in the electronic catalogue, a permutation $p$ means
$N^{\rm census}_{ij}{}^k=N^{\rm source}_{p(i),p(j)}{}^{p(k)}$.
There can be more than one such permutation; the catalogue stores the
lexicographically first one. The six matrices are independently
transcribed in the constructor and all $216$ coefficients are compared.
Schopieray's Theorem~1 proves that this based ring admits no
pseudo-unitary categorification. Its exact existence record is therefore
\[
\begin{array}{c|ccccc}
 &\text{general}&\text{pivotal}&\text{spherical}&\text{pseudo-unitary}&\text{unitary}\\\hline
\RR{2}{124}&\text{yes}&\text{yes}&\text{yes}&\text{no}&\text{no}.
\end{array}
\]
The last entry follows because a unitary fusion category is
pseudo-unitary. This is a ring-level nonexistence theorem, not only the
failure of a particular spherical structure to have positive dimensions.
It does not contribute to the count of \emph{general} exclusions.

\subsubsection{A nilpotent extension and an Izumi even part}
Write $H_8$ for the Kac--Paljutkin Hopf algebra of dimension eight.
Choi--Lu--Sun give extensions of $\Rep(H_8)$ with one additional simple
object in \cite[equation~(1.4), Section~7.1 and Tables~3--4]{CLS}.
With basis $(1,a,b,ab,N,D)$ and $a^2=b^2=1$, their fusion rules are
\begin{equation}\label{eq:h8-extension}
 \begin{gathered}
 gN=N,\quad gD=D,\quad N^2=\sum_{g\in C_2^2}g,\\
 ND=DN=2D,\qquad D^2=\sum_{g\in C_2^2}g+2N.
 \end{gathered}
\end{equation}
The relation $gN=N$ is meant for each $g\in C_2^2$, as is $gD=D$.
The source provides existence; it is not inferred from associativity of
these integer rules.

For a separate categorical justification of unitarizability, the
successive adjoint subcategories determined by
\eqref{eq:h8-extension} have bases
\[
 \{1,a,b,ab,N,D\}\supset\{1,a,b,ab,N\}
 \supset\{1,a,b,ab\}\supset\{1\}.
\]
Thus every categorification with these rules is nilpotent. A nilpotent
fusion category is weakly group-theoretical and is unitarizable by
\cite[Theorem~5.20]{GHR}. The program recomputes this adjoint series by
support closure. The source tensor matches $\RR{2}{69}$ with the identity
permutation. This proves unitary \emph{existence}; neither nilpotency
alone nor a tensor-axiom check is used to assert existence.

The dual even part of the Izumi $C_4$ subfactor is a unitary fusion
category. Its multiplication table in \cite[p.~27]{IMP}, in the source
order $(1,\kappa,\beta\kappa,\chi,\sigma,\beta)$, matches
$\RR{2}{98}$ after the permutation
\[
 (0,5,1,2,4,3).
\]
The constructor recovers every entry of that table before making the
comparison. The equality is with $\RR{2}{98}$, not with the
dimension-matching tensor $\RR{2}{41}$. Equality of dimension lists
would not establish this identification.

\subsubsection{A free central-boson quotient}
A \emph{central boson} here means an invertible object $\widetilde g$ in
the unitary Drinfeld center $Z(\mathcal D)$, generating a symmetric
subcategory equivalent to $\Rep(C_2)$, whose underlying object $g$ in
$\mathcal D$ is nontrivial. This includes the half-braiding and the
Tannakian structure. Centrality of the basis element $[g]$ in
$K_0(\mathcal D)$ is not a substitute for this hypothesis.

\begin{proposition}[Free central-boson quotient]\label{prop:central-boson}
Let $\mathcal D$ be a unitary fusion category and let $\widetilde g$ be
a central boson as above. Suppose tensoring with $g$ has no fixed point
on $\Irr(\mathcal D)$. Then the right modules over the central algebra
$A=1\oplus g$ form a unitary fusion category $\mathcal D_A$. Its simples
are indexed by the two-element $g$-orbits. For orbit representatives
$i,j,k$ its fusion coefficients are
\begin{equation}\label{eq:orbit-quotient}
 \widetilde N_{[i],[j]}^{[k]}=N_{ij}^{k}+N_{ij}^{gk}.
\end{equation}
\end{proposition}
\begin{proof}
The algebra $A$ is the image of the regular function algebra of $C_2$
under the central Tannakian embedding. It is connected and separable,
and is commutative as a central algebra. The half-braiding makes right
$A$-modules into a tensor category with product $\otimes_A$; see
\cite[Theorem~8.23.3]{EGNO}. Use the standard dagger Frobenius structure on the regular algebra.
Work first with its unitary modules. Every algebraic module is a
summand of a free module by separability; an idempotent in the
finite-dimensional $C^*$-endomorphism algebra of a free unitary module
is conjugate to an orthogonal projection. Thus these unitary modules
represent the whole algebraic module category up to equivalence.
Relative tensor products are the images of the self-adjoint
separability projections. The unitary associativity maps of $\mathcal D$ descend
to unitary maps between their images, and the module morphism spaces
inherit the adjoint operation. Thus $\mathcal D_A$ is a unitary fusion
category.

The free-module functor $F(i)=i\otimes A$ is a tensor functor. By
adjunction,
\[
 \operatorname{Hom}_A(F(i),F(j))
 \cong\operatorname{Hom}_{\mathcal D}(i,j\oplus jg).
\]
For simple $i$, freeness of the $g$-action makes the endomorphism
algebra of $F(i)$ one-dimensional, so $F(i)$ is simple. The same formula
shows $F(i)\cong F(j)$ precisely when $j=i$ or $j=gi$. Every simple
module is a summand of a free module, since $A$ is separable, and hence
is one of these simple free modules. Finally
\[
 F(i)\otimes_A F(j)\cong F(i\otimes j).
\]
Collecting the two constituents $k,gk$ in each orbit gives
\eqref{eq:orbit-quotient}. It also proves independence of the chosen
representatives.
\end{proof}

Apply the proposition to
\[
 \mathcal D=\mathrm{Ising}\boxtimes PSU(2)_6,
 \qquad g=(\psi,V_6).
\]
In the principal unitary braidings, $\psi$ and $V_6$ are order-two
fermions. Their product has self-braiding $+1$ and generates a
Tannakian $\Rep(C_2)$, with central lift supplied by the product braiding.
The simple labels of $PSU(2)_6$ are $V_0,V_2,V_4,V_6$; tensoring by
$V_6$ exchanges $V_0,V_6$ and $V_2,V_4$. Therefore the diagonal action
on the twelve simples of $\mathcal D$ is free. Formula
\eqref{eq:orbit-quotient} gives a rank-six unitary category, whose
tensor is $\RR{2}{94}$. The checker reconstructs the parent tensor,
checks all twelve action labels, computes the six orbits and verifies
the displayed formula and its independence of representatives. Its
orbit order is specified in \filename{data/diagonal_fermion_quotient.json};
the final census permutation is $(0,3,1,2,4,5)$.

The construction uses \emph{all} right $A$-modules with the tensor
product supplied by the central lift. It does not identify them with
the local modules of a braided algebra and does not assert that the
quotient inherits the ambient braiding. In particular, no transparency
hypothesis has been silently added or inferred from the fusion ring.

\subsection{A diagonal fermion with a fixed point}\label{sec:fermionic-product}
The free-action formula of Section~\ref{sec:literature-models} does not
apply to a fixed simple object. Its two resolutions can have different
duality depending on the local rotation data. We give a calculation
which resolves the fixed point and determines every fusion coefficient.

We use positive spherical structures in unitary categories. If $f^2=1$,
$x^*=x$, and $t:f\otimes x\longrightarrow x$ is an isomorphism, denote
by $t^{\mathsf T}:x\otimes f\longrightarrow x$ its pivotal transpose,
using the specified self-dual identifications of $f$ and $x$. Equivalently
it is obtained from $t$ by inserting a coevaluation for $x$ on the right
and evaluating the left pair. Associators are suppressed in this notation.
Normalize $t$ so that the action of two copies of $f$ is the identity.

\begin{lemma}[Duality of a resolved fixed point]\label{lem:fixed-rotation}
Let $(f,c)$ be a central boson of order two, and let $V$ be a self-dual
simple fixed by $f$. Suppose $t:fV\to V$ is a normalized action and
\begin{equation}\label{eq:dual-action-phase}
 t^{\mathsf T}c_{f,V}=\delta t,\qquad \delta\in\{1,-1\}.
\end{equation}
For the central algebra $A=1\oplus f$, the two simple left $A$-modules
with underlying object $V$ have actions $t$ and $-t$. Their duality is
$V_\epsilon^*=V_{\delta\epsilon}$, where $\epsilon=\pm1$.
\end{lemma}
\begin{proof}
The two actions satisfy the module equation because $t$ is normalized.
They are simple and inequivalent since $V$ is simple and a scalar
intertwiner cannot change the sign of its nonzero action.
Their sum is the free module on $V$: the commuting algebra of that free
module is the two-dimensional semisimple algebra of its stabilizer.

The dual left action is obtained by transposing the given action and
then moving $f$ from the left to the right of $V^*$ with its specified
half-braiding. This is exactly the composite in
\eqref{eq:dual-action-phase}. The normalization $f^2=1$ makes the inverse
action equal to the action, and the symmetric Frobenius structure of
$A$ supplies the ordinary pivotal transpose. Thus its sign is
$\delta\epsilon$. Alternatively, on the one-dimensional ordinary
vacuum channel of $V\otimes V$, the projector for the relative tensor
product of signs $\epsilon,\eta$ has scalar
$(1+\epsilon\eta\delta)/2$. The unit occurs precisely when
$\eta=\delta\epsilon$, giving the same duality rule.
\end{proof}

\begin{theorem}[Equal-phase fermionic product]\label{thm:fermionic-product}
Let $\mathcal E_m$ and $\mathcal E_n$ be unitary fusion categories with
respective self-dual bases $1,f,x$ and $1,f',y$ and rules
\[
 f^2=1,\quad fx=xf=x,\quad x^2=1+f+mx,
\qquad
 (f')^2=1,\quad f'y=yf'=y,\quad y^2=1+f'+ny.
\]
Assume specified central fermions $f,f'$ and normalized actions $t,s$
satisfy
\begin{equation}\label{eq:two-fermion-phases}
 t^{\mathsf T}c_{f,x}=-it,\qquad
 s^{\mathsf T}c_{f',y}=-is.
\end{equation}
The categories have specified positive pivotal identifications used in
these equations. Then $m,n$ are even, and the category of modules over
the central algebra
\[
 A=1\oplus(f\boxtimes f')\quad\text{in}\quad
 \mathcal E_m\boxtimes\mathcal E_n
\]
is a unitary fusion category of rank six. With basis $1,g,a,b,u,v$,
its duality fixes $1,g,a,b$ and interchanges $u,v$. Its complete
multiplication is determined by commutativity and
\begin{gather}
 g^2=1,\quad ga=a,\quad gb=b,\quad gu=v,\quad gv=u,
 \label{eq:ferm-product-1}\\
 a^2=1+g+ma,\quad b^2=1+g+nb,\quad ab=u+v,
 \label{eq:ferm-product-2}\\
 au=av=b+\tfrac m2(u+v),\qquad bu=bv=a+\tfrac n2(u+v),
 \label{eq:ferm-product-3}\\
 \begin{aligned}
 u^2=v^2&=g+\tfrac m2a+\tfrac n2b+\tfrac{mn}{4}(u+v),\\
 uv&=1+\tfrac m2a+\tfrac n2b+\tfrac{mn}{4}(u+v).
 \end{aligned}\label{eq:ferm-product-4}
\end{gather}
If $d_m=\FPdim(x)$ and $d_n=\FPdim(y)$, its dimensions are
$1,1,d_m,d_n,d_md_n/2,d_md_n/2$.
\end{theorem}
\begin{proof}
The product of the specified fermions is a central boson: its
self-braiding is $(-1)(-1)=1$. Its central regular algebra $A$ is connected
and separable. Modules over this algebra have relative tensor product,
with tensor unit $A$, and form a unitary fusion category. This uses all
modules over a \emph{central} algebra; no braiding on the quotient is
asserted. Write $Q$ for the tensor free-module functor.

There are four free orbits and one fixed point in the nine simples of
the Deligne product. The corresponding modules are
\[
 1=Q(1,1),\quad g=Q(f,1),\quad a=Q(x,1),\quad b=Q(1,y),
 \qquad Q(x,y)=u\oplus v.
\]
On the fixed point $V=x\boxtimes y$, use the action $t\boxtimes s$.
Its square is one. The transpose and half-braiding in a Deligne product
factor, so \eqref{eq:two-fermion-phases} gives
\[
 (t\boxtimes s)^{\mathsf T}c_{f\boxtimes f',V}
   =(-i)^2(t\boxtimes s)=-(t\boxtimes s).
\]
Lemma~\ref{lem:fixed-rotation} proves $u^*=v$. The other four objects
are self-dual because they are tensor free modules on self-dual objects.
The module dimension equals the underlying dimension divided by
$\dim A=2$, giving the asserted dimensions and rank.

Tensoring a resolved fixed point by $g$ changes its action sign.
Indeed, to form the relative product with $Q(f,1)$ the diagonal algebra
line is moved past $(f,1)$; the scalar is $c_{f,f}=-1$. Identifying
$(f,1)V$ with $V$ by $t\boxtimes\mathrm{id}$ leaves no other scalar,
by the normalized square of $t$. Hence $gu=v$, $gv=u$.
Taking duals gives $ug=v$, $vg=u$, so $g$ is central on these objects.
The free-module tensor identity gives $g^2=1$, $ga=ag=a$, $gb=bg=b$
and \eqref{eq:ferm-product-2}.

Since $g$ fixes $a$, the two products $au,av$ are equal, and likewise
$ua=va$. Their respective sums, computed by the tensor functor $Q$,
both equal $2b+m(u+v)$. This proves the first half of
\eqref{eq:ferm-product-3} on both sides and also forces $m$ even.
The second half and the evenness of $n$ follow similarly. Thus all
products involving $a,b,g$ commute. Since $v=gu$ and $g$ is central,
$u$ and $v$ commute as well.

Now $u^2=v^2$ and $(u^2)^*=v^2=u^2$. The coefficient of $1$ in $u^2$
is zero and that of $g$ is one. Frobenius reciprocity and
\eqref{eq:ferm-product-3} give coefficients $m/2,n/2$ on $a,b$.
Self-duality of $u^2$ makes the coefficients of $u,v$ equal. Put
$e=d_md_n/2$ and call the common coefficient $k$. From
$d_m^2=2+md_m$ and $d_n^2=2+nd_n$ one obtains
\[
 e^2=1+\tfrac m2d_m+\tfrac n2d_n+\tfrac{mn}{2}e.
\]
Dimension comparison gives $k=mn/4$ and proves the square formula.
Multiplication by $g$ then gives the formula for $uv$.
\end{proof}

\begin{corollary}[Three tensor identifications]\label{cor:fermionic-models}
The cases $(m,n)=(0,0),(0,2),(2,2)$ of
Theorem~\ref{thm:fermionic-product} have unitary realizations and identify,
respectively, the census rings
\[
 \RR{1}{3},\qquad \RR{2}{14},\qquad \RR{2}{40}.
\]
In particular, $\RR{2}{40}$ is the diagonal-fermion quotient of two
\emph{identical} unitary half-$E_6$ categories.
\end{corollary}
\begin{proof}
For $m=0$ use the Ising category with positive dimension $\sqrt2$ and
Frobenius--Schur sign $+1$. In the conventions of
\cite[Section~2.1, equations~(9)--(13)]{ALW}, choose $A=e^{3\pi i/8}$;
then $d=-A^2-A^{-2}=\sqrt2$ and $A^4=-i$. The displayed fermion
half-braiding and pivotal normalizations give
\eqref{eq:two-fermion-phases}.

For $m=2$ use unitary half-$E_6$, the category with
\eqref{eq:E6}. The local data in
\cite[Section~7.1, equation~(235), Section~7.2, equation~(240),
Appendix~D.1, equation~(467)]{ALW} give the normalized action square one,
pivotal rotation scalar one, and the half-braiding scalar $-i$ on the
noninvertible object. They therefore give exactly the second input
needed in \eqref{eq:two-fermion-phases}. Categorical existence and
unitarity are those of the Ising and half-$E_6$ constructions, not
inferences from these few scalar identities alone.

The code reconstructs \eqref{eq:ferm-product-1}--\eqref{eq:ferm-product-4}
for the three pairs, checks every fusion-ring axiom, and matches all
$216$ coefficients under unit-fixing permutations. For $(2,2)$ the
census basis is $(1,u,v,g,a,b)$, giving the stated identification.
\end{proof}

The two half-$E_6$ factors must not be replaced without comment by a
category and its complex conjugate. Conjugation changes $-i$ to $+i$;
the product of the rotation phases is then $+1$, and
Lemma~\ref{lem:fixed-rotation} makes the two resolutions self-dual.
The relative-product description of the $3311$ category in
\cite[Section~5.1]{IJMS} uses the opposite/conjugate factors.
Our non-self-dual resolution is therefore not a dimension-only
reidentification of that model. The construction uses central boson
condensation with a fixed point, not the earlier free-orbit formula.

\subsection{The exceptional $G_2$ level-four model}\label{sec:g2-e4}
The cyclic-four tensor $\RR{2}{4}$ has a unitary realization.
We identify it with the exceptional
quantum subgroup $\mathcal E_4(G_2)$ associated with the conformal
embedding
\begin{equation}\label{eq:g2-embedding}
 G_{2,4}\ \subset\ \operatorname{Spin}(14)_1.
\end{equation}
The embedding, vacuum algebra, annular matrix and branching matrix
used below are given in \cite[Section~4.5, pp.~29--32]{CRT}. We use
existence of this conformal-embedding algebra, not existence inferred
merely from a commuting modular-invariant matrix.

Let $\mathcal B$ be the unitary WZW category of $G_2$ at level four,
with the long fundamental weight first, as in \cite{CRT}. The vacuum
algebra is
\[
 A=V_{(0,0)}\oplus V_{(0,3)},\qquad
 \dim A=6+2\sqrt6.
\]
It is the connected commutative separable algebra, with its unitary
structure, supplied by \eqref{eq:g2-embedding}. Thus $\mathcal B_A$,
the category of all right modules with tensor product over $A$, is an
ordinary unitary fusion category. Its subcategory of local modules is
the level-one category of $\operatorname{Spin}(14)$.
This distinction between all modules and local modules matters: the
former has six simples, the latter only four.

\begin{proposition}\label{prop:g2-e4}
The Grothendieck ring of $\mathcal B_A$ is based-isomorphic to
$\RR{2}{4}$. In particular this ring has a unitary categorification.
\end{proposition}
\begin{proof}
In the source ordering $(1,X,Y,h,g,g^{-1})$, the basic annular matrix
and branching matrix are
\[
 N_X=\begin{pmatrix}
 0&1&0&0&0&0\\1&2&2&1&0&0\\0&2&2&0&1&1\\
 0&1&0&0&0&0\\0&0&1&0&0&0\\0&0&1&0&0&0
 \end{pmatrix},\qquad
 E=\begin{pmatrix}
 1&0&0&0&0&0\\0&1&0&0&0&0\\0&1&1&0&0&0\\
 0&0&1&1&0&0\\1&1&1&0&0&0\\0&1&1&0&1&1\\
 0&1&0&1&0&0\\0&1&1&0&0&0\\0&0&1&0&0&0
 \end{pmatrix}.
\]
The source weights, indexing the rows of $E$, are
\[
 (0,0),(0,1),(0,2),(1,0),(0,3),(1,1),(0,4),(1,2),(2,0).
\]
The free-module image of $V_{(0,1)}$ is the simple $X$, by the second
row of $E$, so its annular action really is multiplication by $X$.
The free-module functor is central, hence $X$ commutes with every
object. The local simples are $1,h,g,g^{-1}$, corresponding to the
vacuum, vector, and two spinor representations. The level-one fusion
group of $D_7$ is $P/Q\simeq C_4$; equivalently its vector is the square
of either spinor. Thus $g^2=h$ and $g^4=1$. These four objects remain
distinct and invertible in $\mathcal B_A$, since the local subcategory
is a full tensor subcategory.

Reading the displayed matrix gives
\[
 X^2=1+h+2X+2Y,\qquad hX=X,\qquad gX=g^{-1}X=Y.
\]
The image $X$ is self-dual, since $V_{(0,1)}$ is self-dual. Moreover
$Y=gX$, so $Y^*=Xg^{-1}=Y$. The category is tensor generated by the
central object $X$ and the cyclic local group, so its ring is
commutative. Associativity now determines all remaining products:
\begin{equation}\label{eq:g2-e4-rules}
 \begin{gathered}
 g^4=1,\quad gX=Y,\quad gY=X,\\
 X^2=Y^2=1+g^2+2X+2Y,\qquad
 XY=g+g^3+2X+2Y.
 \end{gathered}
\end{equation}
These are the required full fusion rules, with census order
$(1,g,g^{-1},X,Y,h)$; in source coordinates the permutation is
$(0,4,5,1,2,3)$. The resulting $216$ coefficients, including the
nontrivial duality $g^*=g^{-1}$ and the self-duality of $X,Y$, agree
with $\RR{2}{4}$.
\end{proof}

The mathematical inputs to this identification and its arithmetic
checks are separated. Existence and unitarity come from the conformal
embedding. The source annular and branching matrices determine the
module action; the cyclic local group then determines its full fusion
ring. Independently, the constructor calculates all nine $G_{2,4}$
fusion matrices by Weyl alternants in $\QQ(\zeta_{24})$, and verifies
\[
 N_i^{\mathcal B}E
   =E\Bigl(\sum_a E_{ia}N_a^{\mathcal B_A}\Bigr),\qquad
 EE^{\mathsf T}=I+N_{(0,3)}^{\mathcal B}.
\]
The latter is free-module adjunction for the stated vacuum algebra.
All entries, all unit-fixing identifications with the census, and the
cyclic group action are exact. A dimension match alone would not prove
the proposition: it would fail to distinguish the cyclic-four ring
from the Klein-four rings with the same dimension list.

The dimensions are $(1,1,1,1,2+\sqrt6,2+\sqrt6)$ and
$\FPdim(\mathcal B_A)=24+8\sqrt6$. The published construction and its
exact tensor identification provide a positive control for the finite
necessary systems in Section~\ref{sec:frontier}.

\subsection{Inversion equivariantization of an odd near-group category}
\label{sec:inversion-model}
The positive case at multiplicity three is an existing unitary model,
not a proposed solution of new pentagon equations. We use the coherent
order-two action constructed in \cite[Section~13.3]{Izumi}. The following
formulation separates that categorical input from the reconstruction of
the entire fusion tensor.

\begin{proposition}[An inversion-equivariantization constructor]
\label{prop:inversion-model}
Let $n$ be odd, and let $\mathcal N$ be a unitary near-group category with
invertibles $\alpha_j$, $j\in\ZZ/n\ZZ$, and one further simple $\rho$,
with
\[
 \alpha_j\rho=\rho\alpha_j=\rho,\qquad
 \rho^2=\bigoplus_j\alpha_j\oplus n\rho.
\]
Assume a coherent unitary $C_2$-action fixes $\rho$, takes $\alpha_j$ to
$\alpha_{-j}$, fixes the unit channel of $\rho^2$, and acts on a chosen
basis $(T_j)_j$ of $\operatorname{Hom}(\rho,\rho^2)$ by $T_j\mapsto T_{-j}$.
Assume moreover that the action on the pointed part has the standard
inversion equivariantization $\Rep(\widehat C_n\rtimes C_2)$.
Then $\mathcal N^{C_2}$ is a unitary fusion category of rank $(n+7)/2$.
Its simples are
\[
 1,g,x_+,x_-,v_1,\ldots,v_{(n-1)/2}.
\]
Writing $v_0=1+g$ and $v_j=v_{-j}$, with indices modulo $n$, its products
are determined by
\[
 g^2=1,\qquad gx_+=x_-,\qquad gx_-=x_+,\qquad gv_i=v_i.
\]
\begin{align*}
 v_i v_j&=v_{i+j}+v_{i-j},\\
 v_i x_\pm=x_\pm v_i&=x_++x_-,\\
 x_+^2=x_-^2
 &=1+\sum_i v_i+\frac{n+1}{2}x_++\frac{n-1}{2}x_-,\\
 x_+x_-=x_-x_+
 &=g+\sum_i v_i+\frac{n-1}{2}x_++\frac{n+1}{2}x_-.
\end{align*}
All its simples are self-dual. The dimensions are $1,1,d,d,2,\ldots,2$,
where $d>0$ satisfies $d^2=n+nd$.
\end{proposition}
\begin{proof}
For an action on a semisimple category, equivariant simple objects are
obtained from orbits of simple objects and irreducible projective
representations of their stabilizers. In the present explicitly given
action, the fixed simples $1,\rho$ have their indicated equivariant
structures and each admits the two sign choices. Denote these by $1,g$
and $x_+,x_-$. Each nontrivial inversion orbit has trivial stabilizer
and yields a single $v_i$ with underlying object
$\alpha_i\oplus\alpha_{-i}$. This proves the rank and dimension statements.
Tensoring with $g$ changes the sign of an equivariant structure and
therefore exchanges $x_+,x_-$ while fixing every $v_i$.

The pointed equivariantization hypothesis gives the displayed products
of the $v_i$. For $v_i\rho$, the two summands of the underlying object
are exchanged by the action. Its $+1$ and $-1$ eigenspaces are both
one-dimensional, proving $v_i x_\pm=x_++x_-$, and likewise for products
in the opposite order.

The unit channel in $\rho^2$ contributes $1$ to $x_+^2$. The other
invertible channels are paired and contribute the $v_i$. In the
$\rho$-multiplicity space, inversion fixes exactly $T_0$ and pairs the
other $n-1$ basis vectors. Its eigenspaces have dimensions $(n+1)/2$
and $(n-1)/2$. This proves the formula for $x_+^2$. Multiplying by $g$
proves the mixed product, and multiplying twice proves $x_-^2=x_+^2$.
The unit coefficients in these formulas show self-duality. A coherent
unitary action has a unitary equivariantization; thus the construction
uses actual associativity data, not merely an automorphism of a based
ring.
\end{proof}

\begin{corollary}\label{cor:inversion-five}
The ring $\RR{3}{218}$ admits a unitary categorification.
\end{corollary}
\begin{proof}
Izumi's $C_5+5$ category has the inversion symmetry of
\cite[Example~13.8]{Izumi}. Section~13.3 constructs an actual order-two
automorphism commuting with $\rho$, not just an outer automorphism
class. Its action on the isometries is $S_j\mapsto S_{-j}$ and
$T_j\mapsto T_{-j}$. Theorems~13.17--13.18 of that source establish
precisely the hypotheses and product decompositions used above,
including the pointed representation subcategory.

The source scalar symmetries can also be checked directly: with
$\zeta=\exp(2\pi i/5)$, the quadratic phases are
$a(j)=\zeta^{2j^2}$ and the nonzero $b$-phases have exponents
$(4,1,1,4)$ at $j=1,2,3,4$, each divided by $\sqrt5$.
Both lists are inversion invariant. The dimension used here is derived
from the fusion equation,
\[
 d=\frac{5+3\sqrt5}{2},\qquad d^2=5+5d,
 \qquad\dim(\mathcal N^{C_2})=45+15\sqrt5.
\]
Thus no independently copied numerical dimension enters the constructor.
The source supplies existence and the coherent action; the finite
checks validate the specified permutations and scalar symmetries.

In the proposition's order $(1,g,x_+,x_-,v_1,v_2)$, the census order is
$(1,g,x_-,x_+,v_1,v_2)$. The function
\filename{inversion_rules(5)} constructs every product and the recorded
permutation $(0,1,3,2,4,5)$ identifies all $216$ structure constants with
$\RR{3}{218}$. This is a based-ring identification, not a match of
Frobenius--Perron dimensions alone.
\end{proof}

The generic formula is not a claim that the categorical input exists
for every odd $n$. The positive registry supplies only the cited $n=5$
input. This constructor and its full based isomorphism are included in
Table~\ref{tab:models} and in the electronic model catalogue.

\subsection{Higher subfactor and affine-group models}\label{sec:through7-models}
\subsubsection{The six-object Asaeda--Haagerup category}
The unitary category $\mathcal{AH}_1$ has the six-object fusion table
in \cite[Table~1]{GSAH}. In its source order
$(1,\psi,\chi,\sigma,\zeta,\tau)$, the permutation
$(0,1,3,2,5,4)$ identifies all $216$ structure constants with
$\RR{6}{1764}$. The supplement encodes the source table separately from
the census and checks this complete equality. Existence comes from the
published subfactor category, not from the arithmetic of that table.
This is a known unitary model. Its global dimension coinciding with
that of another tensor, in particular $\RR{3}{215}$, proves neither a
based-ring isomorphism nor a Morita equivalence.

\subsubsection{A regular-orbit representation-ring constructor}
\begin{proposition}\label{prop:regular-orbit-model}
Let $V$ be a nontrivial finite abelian group and suppose that $H$ acts
freely and transitively on $\widehat V\setminus\{1\}$. Put
$h=|H|=|V|-1$. The irreducibles of $G=V\rtimes H$ consist of the
inflations of $\Irr(H)$ and one additional self-dual representation
$\rho$ of dimension $h$. Their products satisfy
\begin{equation}\label{eq:regular-orbit-products}
 \sigma\rho=(\dim\sigma)\rho,\qquad
 \rho^2=\sum_{\sigma\in\Irr(H)}(\dim\sigma)\sigma+(h-1)\rho.
\end{equation}
Consequently these rules, together with $\Rep(H)$, define a unitary
categorical model, whenever the stated group action is supplied.
\end{proposition}
\begin{proof}
The irreducibles of a semidirect product by an abelian normal subgroup
are obtained from $H$-orbits in $\widehat V$ and representations of their
stabilizers. The trivial orbit gives the inflated representations of
$H$. The single nontrivial orbit has trivial stabilizer and gives
$\rho=\operatorname{Ind}_V^G\lambda$, for any nontrivial $\lambda$.
Inversion preserves this orbit, so $\rho$ is self-dual. Restriction of
an inflated $\sigma$ to $V$ is $(\dim\sigma)1$, and induction gives
the first equation. Its reciprocity then gives the multiplicity
$\dim\sigma$ of each inflation in $\rho^2$. If the remaining
multiplicity of $\rho$ is $a$, comparing dimensions yields
$h^2=\sum_\sigma(\dim\sigma)^2+ah=h+ah$, so $a=h-1$.
The representation category of a finite group is unitary.
\end{proof}

Take $V=\FF_3^2$ and let $H$ be generated by
\[
 I=\begin{pmatrix}0&1\\2&0\end{pmatrix},\qquad
 J=\begin{pmatrix}1&1\\1&2\end{pmatrix}.
\]
The identities $I^2=J^2=-1$ and $IJ=-JI$ identify $H$ with $Q_8$.
Its orbit on the eight nonzero vectors is free and transitive; the
contragredient action has the same property. Thus
Proposition~\ref{prop:regular-orbit-model} gives a six-object category
of dimension $72$, with simple dimensions $(1,1,1,1,2,8)$ and
multiplicity seven. It realizes $\RR{7}{572}$.

The constructor generates the actual $72$-element semidirect product,
checks its multiplication and action, and computes the four linear
characters, the inflated two-dimensional quaternion character, and
the induced eight-dimensional character. It verifies orthogonality,
completeness by the sum of squares, and all tensor coefficients by
character inner products. The resulting tensor is then compared to the
census. This test is an application of the proposition to explicit
group data, not a rule accepting all abstract rings with these dimensions.

\subsection{The complete catalogue and group-theoretical constructions}\label{sec:model-catalogue}
The positive part of the classification is recorded as explicitly as the
negative part. Table~\ref{tab:models} contains one row for every one of the
$\PositiveCount$ positive rings, including all twenty-one rings of
multiplicity one. The electronic catalogue records every constructor used,
not merely one representative: there are $\ConstructorCount$ constructor records,
with an explicit unit-fixing basis permutation from each output to its
census representative. These permutations are verified entry by entry.
The names in the table are presentation labels, not hypotheses in a
categorification test. The symbols U and S distinguish unitary existence
from spherical existence with proved nonexistence of a pseudo-unitary
categorification. The single S entry is not counted among the unitary
models.

\subsubsection{Small constructions and conventions}
Here $PSU(2)_k$ means the even subcategory of the principal unitary
$SU(2)_k$ category, and $D_9$ has order eighteen. The indicated $C_2$ action
on $C_3\times C_3$ is inversion. In $C_3\rtimes C_4$, a generator of $C_4$
acts by inversion on $C_3$; the action therefore factors through $C_2$.
The positive Ising-type rings in the table are denoted $\mathcal I(G,r)$,
where $|G|=4$, a subgroup $H$ of order two is fixed, and $r\in\{0,1\}$.
Their basis consists of $G$ and $\sigma_0,\sigma_1$. An element of
$G/H\cong C_2$ translates the subscripts, and
\[
 \sigma_i\sigma_j=\sum_{a\in G:\,\bar a=i+j+r}a.
\]
The unitary realizations are the Ising zestings constructed in
\cite[Section~4.3]{LPR}; the metaplectic category $SO(5)_2$ and its listed
zest are treated there as well. The program reconstructs these formulas,
ordinary finite-group character products, the Haagerup tensor, and the
$SU(2)$ fusion rules before matching the resulting tensors. Thus the
multiplicity-one positives do not depend on a cardinality-complement
argument.

\begingroup\small\renewcommand{\arraystretch}{1.12}\setlength{\LTcapwidth}{\textwidth}
\setlength{\tabcolsep}{4pt}
\begin{longtable}{@{}lp{8.0cm}cl@{}}
\caption{All positively identified rings in the census. U means a unitary realization; S means a spherical realization and no pseudo-unitary categorification (hence no unitary one). The numbers of types U and S are given in Table~\ref{tab:classification}. Here $R_{m,j}=R^{(6)}_{m,j}$. Every constructor and its based isomorphism are recorded electronically.}\label{tab:models}\\
\toprule Ring & Model or construction & Type & Source\\\midrule\endfirsthead
\multicolumn{4}{l}{\tablename\ \thetable\ continued}\\
\toprule Ring & Model or construction & Type & Source\\\midrule\endhead
\midrule\multicolumn{4}{r}{Continued on the next page}\\\endfoot
\bottomrule\endlastfoot
$R_{1,1}$ & $\mathcal I(C_4,1)$: cyclic, non-self-dual Ising zest & U & \cite{LPR}\\
$R_{1,3}$ & $\mathcal I(C_4,0)$: cyclic, self-dual Ising zest & U & \cite{LPR}\\
$R_{1,4}$ & $\Rep(C_3\rtimes C_4)$ & U & \cite{LPR}\\
$R_{1,6}$ & $\Vec_{C_2}\boxtimes\Vec_{C_3}\simeq\Vec_{C_6}$ & U & \cite{LPR}\\
$R_{1,7}$ & Tambara--Yamagami category for $C_5$ & U & \cite{LPR}\\
$R_{1,9}$ & $\mathcal I(C_2^2,1)$: non-self-dual Ising zest & U & \cite{LPR}\\
$R_{1,10}$ & Zested $SO(5)_2$ category & U & \cite{LPR}\\
$R_{1,11}$ & $\mathrm{Fib}\boxtimes\Vec_{C_3}$ & U & \cite{LPR}\\
$R_{1,14}$ & $\Vec_{S_3}$ & U & \cite{LPR}\\
$R_{1,15}$ & Haagerup category $\mathcal H_3$ & U & \cite{GS}\\
$R_{1,20}$ & $\Vec_{C_2}\boxtimes\mathrm{Ising}$; also $R(0,0)$ & U & \cite{LPR,MY}\\
$R_{1,21}$ & $\Vec_{C_2}\boxtimes\Rep(S_3)$; also $R(0,1)$ & U & \cite{LPR,MY}\\
$R_{1,23}$ & $\Vec_{C_2}\boxtimes PSU(2)_5$ & U & \cite{LPR}\\
$R_{1,24}$ & $PSU(2)_{10}$ (the other $3311$ even-part ring) & U & \cite{LPR,IJMS}\\
$R_{1,26}$ & $SO(5)_2$ & U & \cite{LPR}\\
$R_{1,30}$ & $\Rep(D_9)$, with $|D_9|=18$ & U & \cite{LPR}\\
$R_{1,32}$ & $\Rep((C_3\times C_3)\rtimes C_2)$, inversion action & U & \cite{LPR}\\
$R_{1,34}$ & $PSU(2)_{11}$ & U & \cite{LPR}\\
$R_{1,36}$ & $\mathrm{Fib}\boxtimes\mathrm{Ising}$ & U & \cite{LPR}\\
$R_{1,37}$ & $\mathrm{Fib}\boxtimes\Rep(S_3)$ & U & \cite{LPR}\\
$R_{1,38}$ & $\mathrm{Fib}\boxtimes PSU(2)_5$ & U & \cite{LPR}\\
\addlinespace
$R_{2,2}$ & Fermionic $C_2$-zest of $\operatorname{Ad}(E_6)\boxtimes\Vec_{C_2}$ & U & \cite{DGP}\\
$R_{2,4}$ & $\mathcal E_4(G_2)$: modules for $G_{2,4}\subset\operatorname{Spin}(14)_1$ & U & \cite{CRT}\\
$R_{2,11}$ & $\Rep((C_3\times C_3)\rtimes C_4)$ & U & \cite{EGNO}\\
$R_{2,14}$ & Fermionic grading zest of the $E_6$ category & U & \cite{DGP}\\
$R_{2,34}$ & $TT_3$ & U & \cite{EM}\\
$R_{2,36}$ & $\mathcal C(A_4,1,C_2,1)$ & U & \cite{Integral}\\
$R_{2,40}$ & Diagonal-fermion quotient of two identical half-$E_6$ categories & U & \cite{ALW}\\
$R_{2,49}$ & Quadratic category $R^{\prime}(2)$ & U & \cite{MY}\\
$R_{2,53}$ & $\Vec_{C_2}\boxtimes\operatorname{Ad}(E_6)$; also $R(0,2)$ & U & \cite{MY}\\
$R_{2,69}$ & $C_2$-extension of $\Rep(H_8)$ & U & \cite{CLS}\\
$R_{2,71}$ & $\mathcal C(S_4,\omega,C_2^2,\psi)$ of Theorem~\ref{thm:s4-model} & U & \cite{GNgt}\\
$R_{2,75}$ & $E_6$ modules over $SU(2)_{10}$ & U & \cite{KO}\\
$R_{2,94}$ & Diagonal-fermion quotient of $\mathrm{Ising}\boxtimes PSU(2)_6$ & U & \cite{EGNO}\\
$R_{2,97}$ & Dual Goodman--de la Harpe--Jones $3311$ even part & U & \cite{IJMS}\\
$R_{2,98}$ & Dual even part of the Izumi $C_4$ subfactor & U & \cite{IMP}\\
$R_{2,105}$ & $\mathcal C(A_5,1,D_{10},1)$ & U & \cite{Integral}\\
$R_{2,124}$ & Schopieray: $\mathcal C(\mathfrak{so}_5,3/2)_{\mathrm{ad}}$ & S & \cite{SchNonPU}\\
$R_{2,127}$ & $\mathrm{Fib}\boxtimes\operatorname{Ad}(E_6)$ & U & \cite{KO}\\
$R_{2,128}$ & $D_6$ modules over $SU(2)_8$ & U & \cite{KO}\\
$R_{2,140}$ & Local $D_{10}$ modules over $SU(2)_{16}$ & U & \cite{KO}\\
$R_{2,143}$ & Commutative even part of the Izumi $C_5$ subfactor & U & \cite{EG}\\
$R_{2,144}$ & Principal $G_2$ level-three category & U & \cite{QGsurvey}\\
\addlinespace
$R_{3,47}$ & $\Rep(GL(3,\FF_2))\simeq\Rep(PSL(2,7))$ & U & \cite{Integral}\\
$R_{3,141}$ & $\mathcal C(A_5,1,S_3,1)$ & U & \cite{Integral}\\
$R_{3,218}$ & Inversion equivariantization of Izumi's $C_5+5$ near-group category & U & \cite{Izumi}\\
\addlinespace
$R_{4,249}$ & Sign-swap equivariantization of the equal half-$E_6$ quotient & U & \S\ref{sec:sign-swap}\\
$R_{4,322}$ & Diagonal-fermion quotient of half-$E_6\boxtimes PSU(2)_6$ & U & \S\ref{sec:mixed-fermion}\\
\addlinespace
$R_{5,17}$ & Izumi's $C_5+5$ near-group category & U & \cite{Izumi}\\
\addlinespace
$R_{6,345}$ & $\mathcal C_{4,8,4}$, with ring $R(4,6)$ & U & \cite{MY}\\
$R_{6,1712}$ & Relative Fibonacci center $Z_{\mathrm{Fib}}(TT_3)$ & U & \S\ref{sec:relative-fibonacci}\\
$R_{6,1764}$ & Asaeda--Haagerup category $\mathcal{AH}_1$ & U & \cite{GSAH}\\
\addlinespace
$R_{7,572}$ & $\Rep((C_3\times C_3)\rtimes Q_8)$ & U & \cite{EGNO}\\
$R_{7,2274}$ & Extended Haagerup category $\mathcal{EH}_1$ & U & \cite{BMPS,BPEH}\\
\end{longtable}
\endgroup


\subsubsection{Exact Weyl-character reconstruction}
We use the principal unitary categories $\mathcal C(\mathfrak g,k)$
at positive integral level. Their construction and the Weyl-alcove and
character formulas are reviewed in \cite{QGsurvey}; these are categorical
inputs, not consequences of a numerically fitted modular matrix. Normalize
long roots to have squared length two, let $W$ be the Weyl group, and put
$\rho=\sum_i\lambda_i$. For dominant weights in the level-$k$ alcove define
\begin{equation}\label{eq:weyl-model}
 C_{\lambda\mu}=\frac{\sum_{w\in W}\det(w)
 \exp\bigl(-2\pi i(w(\lambda+\rho),\mu+\rho)/(k+h^\vee)\bigr)}
 {\sum_{w\in W}\det(w)
 \exp\bigl(-2\pi i(w\rho,\mu+\rho)/(k+h^\vee)\bigr)}.
\end{equation}
These are the fusion characters. Their products determine the fusion
coefficients by
\[
 C_{\lambda\mu}C_{\nu\mu}=\sum_\tau N_{\lambda\nu}^{\tau}C_{\tau\mu}.
\]
All denominators and the required character-matrix rank are checked in the
cyclotomic number field. Exact Gaussian elimination determines the
coefficients, which are then checked to be nonnegative integers and to
satisfy every character equation, including columns not used to solve the
linear system. This also works for the root-lattice subcategory of type
$A_2$, using a full-rank submatrix of its restricted characters.

For $G_2$ at level three, the short fundamental weight is first, and the
input data are
\[
 \bigl((\lambda_i,\lambda_j)\bigr)=
 \begin{pmatrix}2/3&1\\1&2\end{pmatrix},\qquad h^\vee=4,
 \qquad\Lambda_3=\{(0,0),(1,0),(2,0),(3,0),(0,1),(1,1)\}.
\]
The calculation takes place in $\QQ(\zeta_{21})$, and produces a
rank-six tensor of multiplicity two.
\begin{proposition}\label{prop:g2model}
The tensor obtained from \eqref{eq:weyl-model} for $(\mathfrak g,k)=(G_2,3)$
is based-isomorphic to $\RR{2}{144}$. In particular this ring admits a
unitary modular categorification.
\end{proposition}
\begin{proof}
The Weyl construction supplies the category and the displayed character
formula supplies its fusion rules. The exact calculation just described
is implemented independently of the census in
\filename{software/quantum_rank2.py}. The resulting $216$ coefficients
and a basis permutation are recorded in the model catalogue; the checker
recomputes them and verifies their equality to the census tensor. This is
a full-tensor identification, not an identification by dimension type.
\end{proof}

\subsubsection{The dual $3311$ even part}
The dual even-part identification is
\begin{equation}\label{eq:ghjmodel}
 \RR{2}{97}\cong K_0\bigl((A_{11}^{\mathrm{even}})^*_{E_6^{\mathrm{even}}}\bigr).
\end{equation}
The category on the right is the dual even part of the
Goodman--de la Harpe--Jones $3311$ subfactor and is unitary.
Its complete multiplication table is Figure~2 of
\cite[Section~5.1]{IJMS}. The independent constructor uses the basis
$(1,B,C,C',D,g)$ of that figure; in particular
\[
 C^2=1+g+2C,\quad (C')^2=1+g+2C',\quad CC'=B+D,
 \quad gC=C,\quad gC'=C',\quad gB=D.
\]
The other products are reconstructed from the same published table and
all tensor axioms and the based isomorphism in \eqref{eq:ghjmodel} are
checked. The other $3311$ even-part ring is already $\RR{1}{24}$.
Neither identification implies that another ring with the same
Frobenius--Perron dimensions is categorifiable; in particular it does not
settle the non-self-dual ring $\RR{2}{40}$.

\subsubsection{A three-double-coset realization criterion}
For finite $G$, a subgroup $H$, and unit-modulus normalized cochains
$\omega\in Z^3(G,\CC^\times)$ and $\psi\in C^2(H,\CC^\times)$ with
$\omega|_H=\delta\psi$, write $\mathcal C(G,\omega,H,\psi)$ for the
category of $\psi$-twisted $H$-bimodules in $\Vec_G^\omega$, with the
conventions of \cite{GNgt}. For $u\in G$ put $H^u=H\cap uHu^{-1}$.
Simple objects over $HuH$ correspond to irreducible projective
representations of $H^u$ with factor
\begin{align}\label{eq:stabilizer-factor}
 \psi^u(a,b)={}&\psi(a,b)\psi(u^{-1}b^{-1}u,u^{-1}a^{-1}u)\nonumber\\
 &\quad\cdot\frac{\omega(a,b,u)\omega(a,bu,u^{-1}b^{-1}u)}
 {\omega(abu,u^{-1}b^{-1}u,u^{-1}a^{-1}u)}.
\end{align}
Such a simple has Frobenius--Perron dimension
$[H:H^u]\dim\pi$. These statements, and support of tensor products on
products of double cosets, are \cite[Theorems~4.1 and~5.1]{GNgt}.
We call a factor on an abelian group nondegenerate when its twisted group
algebra has one irreducible representation; equivalently its commutator
bicharacter has trivial radical.

\begin{theorem}[Three-double-coset realization]\label{thm:three-coset}
Suppose $H$ is abelian of order $h=s^2>1$ and
\[
 G=H\sqcup tH\sqcup HuH,\qquad
 N_G(H)=H\sqcup tH,\qquad H^u=\{1\}.
\]
If the factor $\psi^t$ in \eqref{eq:stabilizer-factor} is nondegenerate,
then $\mathcal C(G,\omega,H,\psi)$ is a unitary categorification of the
commutative fusion ring with basis $\widehat H\sqcup\{X,Y\}$ and rules
\begin{equation}\label{eq:three-coset-ring}
\begin{gathered}
 \chi X=X\chi=X,\qquad\chi Y=Y\chi=Y,\qquad XY=YX=sY,\\
 X^2=\sum_{\chi\in\widehat H}\chi,\qquad
 Y^2=\sum_{\chi\in\widehat H}\chi+sX+(h-2)Y.
\end{gathered}
\end{equation}
Its dimension type is $(1^h,s,h)$.
\end{theorem}
\begin{proof}
The component over $H$ is $\Rep(H)$, so it consists of the $h$ invertibles
indexed by $\widehat H$ with their ordinary group multiplication.
Nondegeneracy of $\psi^t$ gives one simple $X$ of dimension $s$ over
$tH$. The trivial stabilizer of $HuH$ gives one simple $Y$ of dimension
$h$. Hence these are all the simple objects.

The normalizer is a subgroup with quotient $N_G(H)/H$ of order two.
Consequently $tH$ is closed under inversion. Its complement $HuH$ in
$G$ is also closed under inversion and is stable under multiplication
on either side by $N_G(H)$. Since $X$ and $Y$ are unique on their
supports, both are self-dual and all the invertibles fix them on both
sides. Support forces $X^2$ into $\Rep(H)$. Frobenius reciprocity gives
$[X^2:\chi]=[\chi X:X]=1$, proving the first square formula. Similarly,
$XY$ and $YX$ are supported on $HuH$, and dimensions give $XY=YX=sY$.
Reciprocity gives $[Y^2:\chi]=1$ and $[Y^2:X]=s$. The remaining
coefficient is therefore determined by
$h^2=h+s^2+h[Y^2:Y]$, namely $[Y^2:Y]=h-2$.

Finally $\Vec_G^\omega$ is unitary for unit-modulus $\omega$. The
admissible twisted regular algebra of $H$, with its normalized regular
inner product, is a connected separable $C^*$-algebra object. Its
bimodule category is unitary. This is the standard unitary form of the
group-theoretical construction, and proves existence rather than merely
associativity of the displayed ring.
\end{proof}

\subsubsection{An explicit twisted $S_4$ model}
The following cochain certificate makes the realization criterion
applicable to a residual that is not realized by the corresponding
untwisted construction.

Realize $G=PGL(2,\FF_3)\cong S_4$ by its action on
$\mathbb P^1(\FF_3)=\{(1,0),(0,1),(1,1),(1,2)\}$.
In each projective class choose the matrix $s(g)$ whose first nonzero
entry in row-major order is one. Define
\[
 s(g)s(k)=(-I)^{\beta(g,k)}s(gk),\qquad
 \alpha(g)=\begin{cases}0,&\det s(g)=1,\\1,&\det s(g)=2,\end{cases}
 \qquad
 \omega(g,k,l)=(-1)^{\alpha(g)\beta(k,l)}.
\]
Here $\beta$ is a normalized $\FF_2$-valued $2$-cocycle from the central
extension $GL(2,\FF_3)$, and $\alpha$ is a homomorphism. The cup-product
identity therefore proves that $\omega$ is a normalized $3$-cocycle.
Take $H=\{1,a,b,c\}$ represented, respectively, by
\[
 I,\quad
 \begin{pmatrix}0&1\\1&0\end{pmatrix},\quad
 \begin{pmatrix}0&1\\2&0\end{pmatrix},\quad
 \begin{pmatrix}1&0\\0&2\end{pmatrix}.
\]
Its action is the subgroup generated by the two disjoint transpositions
on the displayed ordered projective line, so $H\cong C_2^2$.
Put $\psi(x,y)=i^{q(x,y)}$, with $q(1,x)=q(x,1)=0$ and
\begin{equation}\label{eq:s4-cochain}
 \begin{array}{c|rrr}
 q&a&b&c\\\hline
 a&0&2&0\\
 b&3&0&1\\
 c&1&0&1
 \end{array}\quad\text{in }\ZZ/4\ZZ.
\end{equation}
For every $x,y,z\in H$ the admissibility identity is
\[
 q(y,z)+q(x,yz)-q(xy,z)-q(x,y)
 \equiv 2\alpha(x)\beta(y,z)\pmod4.
\]
It follows either by substitution in the nine-entry table or by the
complete $4^3$-triple check in the supplement.

\begin{theorem}\label{thm:s4-model}
For the cochains just defined,
$\mathcal C(S_4,\omega,C_2^2,\psi)$ is a unitary categorification of
$\RR{2}{71}$.
\end{theorem}
\begin{proof}
Choose
\[
 t=\left[\begin{pmatrix}1&1\\1&2\end{pmatrix}\right],\qquad
 u=\left[\begin{pmatrix}0&1\\1&1\end{pmatrix}\right].
\]
Matrix multiplication over $\FF_3$ gives exactly three double cosets,
of sizes $4,4,16$, represented by $1,t,u$. The first two form
$N_G(H)$, and $H^u$ is trivial. Substitution in
\eqref{eq:stabilizer-factor} gives the following exponent matrix for
$\psi^t$, in the order $(1,a,b,c)$:
\[
 \begin{pmatrix}
 0&0&0&0\\0&1&1&0\\0&3&0&1\\0&2&3&1
 \end{pmatrix}\pmod4.
\]
Its commutator bicharacter has trivial radical: each nonidentity row
has an entry different from the corresponding transposed entry.
Thus its twisted group algebra is $M_2(\CC)$, and
Theorem~\ref{thm:three-coset} applies with $h=4$, $s=2$.

The constructor regenerates the twenty-four projective matrices, their
faithful permutation action, the subgroup, all double cosets, and every
stabilizer factor. It checks the admissibility identity and, in addition
to the cup-product proof, all $24^4$ cocycle identities. It then uses
\eqref{eq:three-coset-ring} to reconstruct the full tensor and checks its
based isomorphism with $\RR{2}{71}$. The complete data are in
\filename{data/twisted_s4_certificate.json}. Thus no inference from
Frobenius--Perron dimensions alone is used.
\end{proof}

The three additions in this section are positive determinations, not
new exclusions. In particular, searching only untwisted group-theoretical
models would miss the last construction. All allocated negative
certificates remain applications of the unified obstruction framework
stated earlier in the paper.



\subsection{Exceptional and near-group source identifications}\label{sec:source-audit}
Two source tensors must be included independently of any obstruction search.
The principal even part $\mathcal{EH}_1$ of the extended Haagerup subfactor
is unitary by its subfactor construction \cite{BMPS}. Its six fusion matrices
are explicitly displayed in \cite[Section~5.1]{BPEH}. In that source order,
the permutation $(0,1,3,4,5,2)$ gives $\RR{7}{2274}$. The largest coefficient
is seven. The constructor transcribes the published matrices and checks
all $216$ coefficients, not merely the dimension field or formal codegrees.

The unitary $C_5+5$ near-group category in \cite[Example~13.8]{Izumi} has
basis $\alpha_0,\ldots,\alpha_4,\rho$ and products
\[
 \alpha_i\alpha_j=\alpha_{i+j},\qquad
 \alpha_i\rho=\rho\alpha_i=\rho,\qquad
 \rho^2=\sum_{i=0}^4\alpha_i+5\rho.
\]
Its tensor identifies $\RR{5}{17}$ with permutation $(0,1,4,3,2,5)$.
The positive root of the dimension equation is $(5+3\sqrt5)/2$.
This is the near-group input itself, distinct from the inversion
equivariantization of Corollary~\ref{cor:inversion-five}. Both source
identifications are independent of the census and both are matched under
all unit-preserving permutations.

\subsection{A mixed diagonal-fermion quotient}\label{sec:mixed-fermion}
Let $\mathcal E$ be the unitary half-$E_6$ category of \eqref{eq:E6}, and
let $\mathcal P=PSU(2)_6$. The object $h$ of $\mathcal E$ has the central
fermion of Lemma~\ref{lem:fermion}. In the standard unitary $SU(2)_6$
category, $V_6$ is an order-two fermion; restricting its half-braiding gives
a central fermion in $\mathcal P$. Thus $(h,V_6)$ is a central boson in
$\mathcal E\boxtimes\mathcal P$.

\begin{proposition}\label{prop:mixed-fermion}
The category of modules over $1\oplus(h,V_6)$ in
$\mathcal E\boxtimes\mathcal P$ is a unitary fusion category with
Grothendieck ring $\RR{4}{322}$.
\end{proposition}
\begin{proof}
The second factor has basis $(V_0,V_2,V_4,V_6)$, with $V_6$ interchanging
$V_0,V_6$ and $V_2,V_4$. The diagonal action is therefore free, even on
objects whose first component is $X$. Proposition~\ref{prop:central-boson}
applies without a fixed-point choice. In the product order
$(1,h,X)\times(V_0,V_2,V_4,V_6)$, its permutation of the twelve simples is
\[
 (7,6,5,4,3,2,1,0,11,10,9,8).
\]
Take orbit representatives $(0,1,2,3,8,9)$. The quotient formula
of Proposition~\ref{prop:central-boson}, applied to the independently
constructed half-$E_6$ and $SU(2)_6$ tensors, gives the census tensor after
permutation $(0,3,1,2,4,5)$. The implementation checks every product and
its independence of the chosen representative. Existence and unitarity
come from the central separable algebra; the tensor comparison identifies
its Grothendieck ring.
\end{proof}

\subsection{A sign-twisted swap on a fermionic product}\label{sec:sign-swap}
The next construction has a categorical hypothesis. Associativity of the
integer tensor for a formal parameter does not provide its input category.

\begin{theorem}[Sign-swap equivariantization]\label{thm:sign-swap}
Let $\mathcal E$ be a unitary fusion category with simples $1,f,x$, a
specified central fermion $f$, and rules
\[
 f^2=1,\quad fx=xf=x,\quad x^2=1+f+mx.
\]
Suppose $d=\FPdim(x)$ is irrational. There is a unitary fusion category
with invertible group $G=C_2\times C_2$ and two self-dual simples $X,Y$,
fixed by all invertibles, with
\begin{align}
 X^2&=\sum_{g\in G}g+mX+2Y,\label{eq:sign-swap}\\
 XY=YX&=2X+2mY,\notag\\
 Y^2&=\sum_{g\in G}g+2mX+m^2Y.\notag
\end{align}
Their dimensions are $2d$ and $d^2$. The Ising input $m=0$ gives
$\RR{2}{69}$, and the half-$E_6$ input $m=2$ gives $\RR{4}{249}$.
\end{theorem}
\begin{proof}
In $\mathcal E\boxtimes\mathcal E$ put $b=f\boxtimes f$ and $A=1\oplus b$.
This is the regular algebra of a central Tannakian subcategory. The
factor-swap tensor functor $S$ preserves the central lift of $b$.
Equip $A$ with the algebra isomorphism $\phi:S(A)\to A$ that is $+1$ on
$1$ and $-1$ on $b$. It satisfies $\phi S(\phi)=\operatorname{id}$.
Transport of module actions through $S$ and $\phi$ defines a coherent
unitary $C_2$ action $T$ on $\mathcal Q=(\mathcal E\boxtimes\mathcal E)_A$.
The category we use is the ordinary equivariantization $\mathcal Q^{C_2}$.

The free $b$-orbits give four simples of $\mathcal Q$: the unit,
$g=(f\boxtimes1)A$, and $a=(x\boxtimes1)A$, $a'=(1\boxtimes x)A$.
The fixed object $x\boxtimes x$ gives two simples $u_+,u_-$, each of
dimension $d^2/2$. To see the action explicitly, choose the right
$f$-action $t:x\otimes f\to x$ with square one, using the module
associativity constraint. Such a normalization exists because
$1\oplus f$ is the regular separable algebra and its two actions on $x$
differ by sign. The right $b$-actions on the fixed object are
$\pm(t\boxtimes t)$. Swap preserves this tensor, while $\phi$ changes its
sign. Hence $T$ interchanges $u_+,u_-$; it also interchanges $a,a'$ and
fixes $1,g$.

An orbit of size two has trivial stabilizer and yields one simple in the
equivariantization; a fixed simple has two equivariant structures, since
projective representations of $C_2$ over $\CC$ can be linearized. Thus
there are four invertibles, one simple $X$ over $a\oplus a'$ and one
simple $Y$ over $u_+\oplus u_-$. Their dimensions are $2d,d^2$.
Their supports are invariant under duality, so $X,Y$ are self-dual.
They are fixed by all invertibles: tensoring preserves their distinct
dimensions and there is only one simple of either dimension.

We check the group law on the four invertibles; dimensions alone would
not distinguish $C_4$ from $C_2\times C_2$. Work in the pointed part
$\mathrm{sVec}\boxtimes\mathrm{sVec}$, with trivial associator and the
specified self-braidings $-1$. Let $e_0,e_1$ be the two homogeneous
components of $g=(f,1)A$, in degrees $(f,1),(1,f)$. An isomorphism
$g\otimes_A g\to A$ has the normalized coefficients
\[
 \mu_{00}=1,\qquad \mu_{01}=b,\qquad
 \mu_{10}=-b,\qquad \mu_{11}=-1.
\]
The minus sign in balancing is forced by the half-braiding of $b$ with
$(f,1)$; it is not a choice of a ring sign. Give $g$ the equivariance
$T(e_0)\mapsto i e_1$, $T(e_1)\mapsto-i e_0$. Its square is the identity.
On each of the four pairs $(e_r,e_s)$, the displayed $\mu$ intertwines
this tensor-square equivariance with $\phi$, which acts on $A$ as
$(1,-1)$. Therefore this lift of $g$ squares to the tensor unit.
Together with the independent sign representation of $C_2$, it generates
$C_2\times C_2$. The six scalar equalities (two for coherence, four for
multiplication) are also checked exactly over $\QQ(i)$.

Frobenius reciprocity now puts each invertible once in $X^2,Y^2$ and none
in $XY$. Write $d^2=2+md$. Comparing coefficients of $1,d$ in
$(2d)^2$, $(2d)d^2$ and $d^4$ forces respectively the coefficients in
\eqref{eq:sign-swap}; their uniqueness uses the irrationality of $d$.
For example $X^2=\sum_g g+aX+cY$ gives $8+4md=4+2c+(2a+cm)d$,
so $c=2$ and $a=m$. For $XY=aX+cY$, comparison gives $c=2m,a=2$;
for $Y^2=\sum_g g+aX+cY$, it gives $c=m^2,a=2m$.
All operations used preserve unitarity.

The Ising and half-$E_6$ inputs and their central fermions are supplied by
standard Ising theory and Lemma~\ref{lem:fermion} (see also \cite{ALW}).
Their tensors, constructed from \eqref{eq:sign-swap}, match the census
with permutations $(0,1,2,3,5,4)$ and the identity. These two categorical
inputs, not arbitrary values of $m$, are admitted by the model registry.
\end{proof}

\subsection{Reconstruction from a relative center}\label{sec:relative-fibonacci}
For a full fusion subcategory $\mathcal D\subset\mathcal C$, the relative
center $Z_{\mathcal D}(\mathcal C)$ consists of objects of $\mathcal C$
equipped with a coherent half-braiding against objects of $\mathcal D$.
It is not the M\"uger centralizer in a braided category. It is a fusion
category, is Morita equivalent to $\mathcal D\boxtimes\mathcal C^{\rm rev}$,
and has dimension $\FPdim(\mathcal D)\FPdim(\mathcal C)$
\cite[Remark~2.4]{GNNrel}. For unitary input, its realization as the dual
of the natural unitary module category gives a unitary fusion category.

\begin{proposition}[Forced relative-center reconstruction]\label{prop:relative-reconstruction}
Suppose $\mathcal D=\langle a\rangle\subset\mathcal C$ is a Fibonacci
subcategory, so $a^2=1+a$. In the simple basis of $\mathcal C$, let $L_a,R_a$
be the matrices of left and right multiplication by $a$, and put
\[
 B=I+L_aR_a,\qquad
 L=(e_0,e_a,e_a,e_0+e_a),\qquad Q=B-LL^{\mathsf T}.
\]
Starting from $Q$, whenever a diagonal entry is one, its column is a
forced nonnegative integral Gram column. Subtract its outer product and
repeat. If this process reaches zero, it determines the entire restriction
matrix $F=(L,H)$ of $Z_{\mathcal D}(\mathcal C)$.

Assume additionally that the high columns are pairwise distinct, distinct
from the low columns, and individually invariant under duality in
$\mathcal C$. Then the relative center is self-dual. Impose the known
$\mathrm{Fib}\boxtimes\mathrm{Fib}^{\rm rev}$ multiplication on its first
four simples and the ring-homomorphism equations for $F$. If these linear
equations in the symmetric nonunit structure constants have a unique
solution, that solution is its complete Grothendieck tensor.
\end{proposition}
\begin{proof}
For the forgetful functor and its adjoint, relative central induction has
underlying object $FI(V)=V\oplus aVa$; equivalently it is the finite
semisimple coend $\bigoplus_{U\in\Irr\mathcal D}U\otimes V\otimes U^*$.
Adjunction gives the Gram identity $FF^{\mathsf T}=B$.
The full inclusion $\mathcal D\subset\mathcal C$ induces a fully faithful
tensor inclusion $Z(\mathcal D)\subset Z_{\mathcal D}(\mathcal C)$:
the half-braiding condition and all morphism spaces between objects of
$\mathcal D$ are unchanged. The standard nondegenerate Fibonacci braiding
identifies $Z(\mathcal D)$ with
$\mathrm{Fib}\boxtimes\mathrm{Fib}^{\rm rev}$, whose four restriction
columns are exactly $L$.

If $Q_{ii}=1$, exactly one remaining simple has restriction coefficient
one in row $i$, and all the others have coefficient zero. Its full
restriction column is therefore $Q_{\bullet i}$. Subtracting its outer
product is mandatory. At zero residual no further nonzero restriction
column is possible. This proves the asserted uniqueness without a
numerical factorization or a conjectural central rank.

The low four simples are self-dual. The additional hypotheses force every
high simple to be self-dual, because duality permutes restriction columns
and no other simple has the same column. A based ring all of whose simples
are self-dual is commutative. Frobenius reciprocity makes its nonunit
coefficients symmetric in all three indices. Its forgetful functor is a
ring homomorphism and the low multiplication is already known. Hence its
actual tensor solves the stated linear system. Uniqueness determines
that tensor; categorical existence was supplied before solving the
system. If a pivot or uniqueness test fails, the recipe is inconclusive,
not an exclusion or a claimed construction from arbitrary ring data.
\end{proof}

\begin{corollary}\label{cor:tt3-relative}
For the unitary category $TT_3$ and either of its Fibonacci generators,
$Z_{\mathrm{Fib}}(TT_3)$ has Grothendieck ring $\RR{6}{1712}$.
\end{corollary}
\begin{proof}
Use the source basis $(1,f^{(2)},\rho,\sigma,\bar\sigma,\mu)$ of
\cite[Definition~2.5]{EM} and $a=\rho$. Its dimension vector is
$(1,\varphi^3,\varphi,\varphi^2,\varphi^2,\varphi)$,
where $\varphi=(1+\sqrt5)/2$. The source provides the unitary category and
its complete multiplication table. The Gram procedure of
Proposition~\ref{prop:relative-reconstruction} gives
\[
 F=\begin{pmatrix}
 1&0&0&1&0&0\\
 0&0&0&0&1&1\\
 0&1&1&1&0&0\\
 0&0&0&0&0&1\\
 0&0&0&0&0&1\\
 0&0&0&0&1&0
 \end{pmatrix}.
\]
Indeed, after removal of the low columns the row of $\sigma$ is a
diagonal-one pivot; after its forced subtraction the row of $f^{(2)}$
is another such pivot. Their removal leaves zero. Denote the columns by
$(1,p,q,pq,X,Y)$. Then $p^2=1+p$, $q^2=1+q$, and the complete remaining
rules are
\begin{align*}
 pX=qX&=Y,&pY=qY&=X+Y,\\
 (pq)X&=X+Y,&(pq)Y&=X+2Y,\\
 X^2&=1+pq+2X+2Y,\\
 XY&=p+q+pq+2X+4Y,\\
 Y^2&=1+p+q+2pq+4X+6Y.
\end{align*}
For detail on uniqueness, there are $\binom{7}{3}=35$ symmetric nonunit
unknowns. The low products and all $36$ restriction-product identities
give $312$ scalar linear equations of exact rank $35$. The verifier
reconstructs this integer system and its unique nonnegative integral
solution, and independently checks every based-ring axiom. The displayed
rules also permit direct substitution into every identity.

The full tensor matches $\RR{6}{1712}$ under $(0,1,2,4,5,3)$.
Repeating the calculation with $a=\mu$ gives a based-isomorphic result.
Applying the same recipe to $\mathcal C=\mathcal D=\mathrm{Fib}$ recovers
$\mathrm{Fib}\boxtimes\mathrm{Fib}^{\rm rev}$ as an independent positive
control. Thus neither the rank-six answer nor its census identification
is an input to the reconstruction.
\end{proof}

\subsection{Three neutral components of conformal extensions}
\label{sec:conformal-models}
The same construction realizes three further census rings. We write
$\mathcal B(\mathfrak g,k)$ for the principal unitary WZW category and
use finite Dynkin labels for its simple objects.

\begin{theorem}\label{thm:conformal-models}
The rings $\RR{4}{22}$, $\RR{4}{91}$ and $\RR{4}{629}$ admit unitary
categorifications. They are the neutral components of the module
categories for the conformal embeddings
\[
 (B_2)_7\subset(D_7)_1,\qquad
 (C_3)_4\subset(B_{10})_1,\qquad
 (B_2)_{12}\subset(E_8)_1,
\]
respectively.
\end{theorem}

\paragraph{Common construction and identification principle.}
A unitary conformal extension supplies an actual connected commutative
separable algebra $A$ in its source category $\mathcal B$.
Consequently $\mathcal E=\mathcal B_A$ is a unitary fusion category,
$F:\mathcal B\to\mathcal E$, $U\mapsto U\otimes A$, is a central
 tensor functor, and the local modules are the level-one target category
$\mathcal L$. These facts use the conformal algebra, not just a modular
invariant \cite[Sections~6.1--6.2]{DMNO}.
In each case $A$ belongs to the even part of the source parity grading.
The two homogeneous parts of any module are submodules. Thus simples
have definite parity, relative tensor product respects parity, and odd
free modules make the grading faithful. Its neutral component
$\mathcal C=\mathcal E_0$ is therefore an actual unitary fusion category.

We identify its full based ring from free induction and local modules.
Two standard tests resolve the few possible tensor completions. First,
\cite[Corollary~3.30]{DMNO} and \cite[Proposition~8.23.11]{EGNO} give
\begin{equation}\label{eq:conformal-neutral-center}
 Z(\mathcal E)\simeq\mathcal B\boxtimes\mathcal L^{\rm rev},
 \qquad
 Z(\mathcal C)\simeq
 \bigl(\Rep(C_2)'\subset Z(\mathcal E)\bigr)_{\Rep(C_2)},
\end{equation}
where $\Rep(C_2)$ is the canonical grading subcategory. If a commutative
candidate ring has formal codegree $f$, the actual known center must
contain a simple of dimension $\dim(\mathcal C)/f$
\cite[Theorem~2.13]{Ostrik3}. Second, indicators can be computed directly
from the actual center and its forgetful functor \cite{NS}.
These tests concern the constructed category and do not presume that
archived modular data have been categorified.

\paragraph{The $B_2$ level-seven component.}
The source \cite[Section~3.5, Figure~11(a)]{CRT} gives
\[
 A=(0,0)\oplus(1,6)\oplus(2,2)\oplus(5,0),
 \qquad \mathcal L=\mathcal B(D_7,1).
\]
Parity is the second Dynkin coordinate modulo two. The six even vertices,
in source order $(1,9,10,2,5,6)$, will be denoted $(1,x,y,g,e,f)$.
Here $e=F(1,0)$ and $g=F(7,0)$ are simple, hence central in the fusion
ring, and $e$ is self-dual. The published vector-action graph gives
\[
 E=N_e=\begin{pmatrix}
 0&0&0&0&1&0\\
 0&1&1&0&0&1\\
 0&1&1&0&0&1\\
 0&0&0&0&1&0\\
 1&0&0&1&0&2\\
 0&1&1&0&2&2
 \end{pmatrix}.
\]
The local category has fusion group $C_4$; an odd local spinor $z$
satisfies $z^2=g$ \cite[Theorem~7.1]{Bock}. The odd simple
$s=F(0,1)$ is self-dual and central, and $zs$ is one of $x,y$.
Consequently $(zs)^*=gzs$: duality on this pair is the action of $g$.
Dimensions force $g$ to fix $e,f$ and either fix or interchange $x,y$.
The first row of $E^2$ gives $e^2=1+g+2f$; the matrix
$(E^2-I-N_g)/2$ is integral only for the interchange. Thus $x^*=y$.

Put $G=N_g$, $H=(E^2-I-G)/2$ and $X=N_x$.
The equation $ef=x+y+2e+2f$ and rigidity imply
\[
 X+GX=EH-2E-2H,\qquad X^{\mathsf T}=GX.
\]
Also $X$ commutes with $E,G,H$, and its unit row is the basis vector
of $x$. These are linear equations over $\QQ$ in the $36$ entries of
$X$. Their coefficient matrix has rank $36$. The unique solution is
nonnegative integral and, with $N_y=GX$ and $N_f=H$, gives exactly all
$216$ coefficients of $\RR{4}{22}$. The accompanying checker extracts
the graph directly from the source PDF and reconstructs this linear
system, its rank, and the full tensor comparison.

\paragraph{The $B_2$ level-twelve component.}
Here \cite[Section~3.6, Figure~13]{CRT} gives
\begin{align*}
 A={}&(0,0)\oplus(0,6)\oplus(1,8)\oplus(3,2)\oplus(3,8)\\
 &\oplus2(4,4)\oplus(6,6)\oplus(7,2)\oplus(12,0).
\end{align*}
The target category $\mathcal L$ is $\Vec$, so $A$ is Lagrangian and
$Z(\mathcal E)\simeq\mathcal B(B_2,12)$.
The even graph component has source vertices $(1,3,7,8,4,6)$.
The central simple $x=F(1,0)$ has multiplication matrix of characteristic
polynomial $t^2(t^4-5t^3+10t-5)$. Every eigenvalue has multiplicity at
most two. A noncommutative matrix summand $M_n(\CC)$ of the complexified
fusion ring would contribute at least $n^2\ge4$ to the multiplicity of
an eigenvalue of this central element. The ring is therefore commutative.

The only duality permutations fixing $1,x$ and commuting with $N_x$
are the identity and the interchange of vertices $7,8$. In the latter
case, Frobenius reciprocity and associativity with $x$ force coefficients
$1/2$ and $3/2$, which are impossible. For identity duality the same
linear constraints leave one parameter $t$. In the indicated basis
$(1,x,u,v,w,z)$,
\[
 N_u=\begin{pmatrix}
 0&0&1&0&0&0\\
 0&0&1&1&1&1\\
 1&1&4-t&t-1&1&2\\
 0&1&t-1&3-t&2&2\\
 0&1&1&2&2&2\\
 0&1&2&2&2&3
 \end{pmatrix}.
\]
All other products are determined and associative; integrality and
nonnegativity give $t\in\{1,2,3\}$. The middle value is exactly
$\RR{4}{629}$. The outer values have formal codegrees $3,15$.

The unique nontrivial invertible $J=(12,0)$ in $\mathcal B(B_2,12)$
is a boson. Thus \eqref{eq:conformal-neutral-center} identifies the
actual center as $(\mathcal B(B_2,12)_{\rm ad})_{\langle J\rangle}$.
Among the $49$ even weights, the involution
$(a,b)\mapsto(12-a-b,b)$ has $21$ ordinary orbits and seven fixed
points. Each ordinary orbit gives one simple of the original Weyl
dimension; each fixed point gives two simples of half that dimension.
This determines all $35$ simple dimensions of the center. With
$D=\dim(A)/2$, exact Weyl arithmetic in $\QQ(\zeta_{60})$ shows that
neither $D/3$ nor $D/15$ occurs. The formal-codegree test excludes both
outer completions, proving the required tensor identification.

\paragraph{The $C_3$ level-four component.}
The adjoint conformal branching theorem
\cite[Theorem~3.9 and Proposition~4.1]{CKMP} gives the three local
Ising modules $1,\psi,\sigma$ with underlying source objects
\begin{align*}
 U(1)=A&=000\oplus210\oplus004\oplus121,\\
 U(\psi)&=200\oplus030\oplus022\oplus301,\\
 U(\sigma)&=2(111).
\end{align*}
Their dimensions are $1,1,\sqrt2$ and their twists are
$1,-1,\zeta_{16}^{5}$. The source parity is $a+c$ modulo two;
there are $19$ even weights.
For free modules the exact induction Gram matrix is
\[
 G_{\lambda\mu}=\langle F(\lambda),F(\mu)\rangle
      =\sum_\alpha[A:\alpha]N_{\lambda\alpha}^{\mu}.
\]
The three local columns are known from branching, while
$\tau=F(010)$ has norm one and is orthogonal to them. Subtracting these
four columns leaves $2vv^{\mathsf T}$ with $v\ge0$ integral and
$v_{002}=1$. The row of norm two must be a sum of two distinct unit
vectors. Rank one forces every other residual row to be proportional
to it. Thus there are exactly two further simples $x,y$, with identical
restriction columns, and $F(002)=x+y$.
Dominance of free induction rules out further simples.

The free image has rank five in the six-dimensional complexified fusion
ring. It is central, so the ring is commutative: a noncommutative matrix
block would make the difference between algebra and center dimensions
at least three. The objects $1,\psi,\sigma,\tau$ are self-dual, leaving
only the alternatives $x^*=y$ and self-duality of both $x,y$.
Source fusion, the induction matrix, the Ising products, Frobenius
reciprocity and associativity give a bounded exact completion problem.
Its four or five free integer parameters lie between zero and four,
as follows from coefficient nonnegativity. In the non-self-dual case
there is one tensor, equal in all $216$ entries to $\RR{4}{91}$.
There are three self-dual possibilities, with respectively
\[
 x^2=y^2=1+2y+2\tau,\qquad
 1+x+y+2\tau,\qquad
 1+2x+2\tau.
\]

To exclude the middle possibility, note that a self-dual simple $X$
has a symmetric self-duality on $X\otimes X$. Indeed, if the pivotal
transpose of $f:X\to X^*$ is $\epsilon f$, then the transpose of
$f\otimes f$, under the reversed dual-tensor identification, is
$\epsilon^2(f\otimes f)$. Every self-dual constituent of multiplicity
one therefore has second indicator $+1$. The middle rule would give
$\nu_2(x)=\nu_2(y)=1$.
On the other hand, the known center
$Z(\mathcal E)\simeq\mathcal B(C_3,4)\boxtimes\mathcal L^{\rm rev}$
and the displayed branching give
\[
 \nu_2(F(\mu))=
 \frac1{\dim\mathcal E}\sum_{\lambda,a,\nu}
 \theta_\lambda^2\theta_a^{-2}d_\lambda d_a
 [U(a):\nu]N_{\nu\lambda}^{\mu}.
\]
Exact evaluation in $\QQ(\zeta_{32})$ gives
$\nu_2(F(002))=0$, contradicting $\nu_2(x+y)=2$.

For the outer two completions the formal codegrees include
$8\pm4\sqrt2$. The boson $J=004$ belongs to $A$, so $F(J)=1$;
the canonical grading subcategory in the known center is generated by
$(J,1)$. Formula~\eqref{eq:conformal-neutral-center} gives
\[
 Z(\mathcal C)\simeq
 (\mathcal B(C_3,4)_{\rm ad})_{\langle J\rangle}
 \boxtimes\mathcal L^{\rm rev}.
\]
The action $J(a,b,c)=(b,a,4-a-b-c)$ has eight ordinary orbits and three
fixed weights among the even weights. Fixed-point splitting gives
rank $14$ for the first factor and rank $42$ for the center.
Its dimension list contains neither
$\dim(\mathcal C)/(8+4\sqrt2)$ nor
$\dim(\mathcal C)/(8-4\sqrt2)$, as checked exactly by the Weyl formula.
The formal-codegree test excludes both outer completions. Only the
non-self-dual target remains, completing the proof of
Theorem~\ref{thm:conformal-models}.

The supplement reconstructs the branching from the abelian ideals of
$C_3$, certifies the source fusion coefficients by exact Weyl-character
identities, and checks every finite completion and dimension comparison.
All three arguments establish categorical existence and full tensor
identifications; they do not classify associators on these rings.

